\section{Algorithms and Checks for Credit Capacity}\label{cr:appendix}

This appendix supplies algorithmic details for \cref{sec:credit}. It uses the same fixed-curve, gross-redemption setting. In particular, the reserve $W$ remains committed while borrowers exercise the repayment rights being tested.

\subsection{Constructing the proportional envelope}\label{cr:algorithm}

Order the loans so that $a_1\ge\cdots\ge a_n$, where $a_i=d_i/\ell_i$. Write $Q_j=\sum_{i\le j}\ell_i$, $B_j=\sum_{i\le j}d_i$, and $Q_0=B_0=0$. For $Q_{j-1}\le q\le Q_j$,
\begin{equation}\label{cr:piecewiseH}
H(q)=B_{j-1}+a_j(q-Q_{j-1}).
\end{equation}
To prove the formula, if an allocation retains some collateral from a smaller-ratio loan while a larger-ratio loan is not fully retained, transfer the smaller of the available two amounts to the larger-ratio loan. Total collateral stays fixed and remaining debt does not decrease. Repeating this exchange gives \eqref{cr:piecewiseH}. Ties have no effect on $H$.

Suppose the input represents $p$ by $m$ intervals with rational endpoints, rational affine prices $p(q)=c_j+k_jq$ on each interval, $k_j\ge0$, and non-negative jumps between intervals. The reserve function is then continuous and piecewise quadratic, with exact rational values at rational arguments. Restrict all intervals to $[0,L]$ and merge their endpoints with the $Q_j$. On any interval of this merged partition,
\[
H(q)-R(q)=b+aq-cq-\tfrac12 kq^2
\]
for an appropriate constant $b$. A maximum on that interval is attained at an endpoint or, if $k>0$, at $q=(a-c)/k$ when that point lies in the interval. If $k=0$, endpoints suffice, including the case of a constant difference. Evaluating this finite candidate list therefore finds $W_{\min}$ exactly.

Sorting requires $O(n\log n)$ comparisons. With price intervals already ordered, the merged scan takes $O(n+m)$ arithmetic operations; the number of candidates is linear in $n+m$. Rational additions, products and comparisons have polynomial bit complexity in the input length. This establishes polynomial-time exact evaluation for this representation. It makes no such computational claim for an unspecified real-valued function $p$.

Concavity also gives a useful check: the slopes of $H-R$ cannot increase. A maximizer lies where these slopes change from non-negative to non-positive, or at an endpoint. The candidate method used in the script checks every merged interval, rather than relying on an approximate numerical root. The equal-ratio comparison in \cref{cr:capacity} needs no uniqueness of this maximizer. The closed form \eqref{cr:linearcapacity} follows by minimizing $P+kq/2+W/q$: for $W>0$ its unconstrained minimum is at $q=\sqrt{2W/k}$; if that depth exceeds $L$, the constrained minimum is at $L$. For $W=0$ the infimum is the limit $P$ at zero.

\subsection{Guards for a proportional book}\label{cr:guards}

Let the extended invariant be custody, the amount-domain conditions, and $W\ge W_{\min}$, with proportional release. This implies the immediate coverage condition because $H(L)=D$:
\[
A=R(S)-D+W\ge R(S)-R(L)\ge0.
\]
After a repayment, each surviving loan has the same ratio $a_i$ and a smaller collateral capacity. Its remaining residual configurations form a subset of the earlier ones, so the maximum defining $W_{\min}$ cannot increase. The repayment leaves $W$ unchanged. A gross sale leaves both the loan book and $W$ unchanged. A sale with a smaller total reserve outflow increases $W$. A purchase whose reserve inflow covers its curve cost also increases $W$ and leaves the existing loan book unchanged.

A withdrawal of $w$ changes only $W$ and cash, so its exact safety guard is $w\le W-W_{\min}$. This bound implies that the withdrawal is affordable: $W_{\min}\ge D-R(L)$ and $R(L)\le R(S)$ give
\[
W-W_{\min}\le W-D+R(L)\le A.
\]
For a proposed new loan or top-up, calculate the resulting debt and locks and the total cash outflow. Accept only if the resulting $W'$ satisfies $W'\ge W_{\min}'$. This implies non-negative resulting cash by the preceding immediate-coverage argument. If the outflow is the gross new debt, then $W'=W$; a retained origination fee increases it. The use of resulting quantities is essential: checking only the old envelope is insufficient.

For a reshape with external inflow $c$, the proposed surplus is $W'=A+c+D-R_{p'}(S)$. Require $W'\ge W_{\min}(p',\{\ell_i,d_i\})$. Proportional rights do not change with the schedule, so this is precisely the condition for the continuation with the new curve. A floor-monotonicity rule can additionally protect the promised floor; it is not a substitute for the new buffer test. For a more general release contract, the test must use its actual residual-pair sets, not the proportional envelope. Nothing here authorizes changing a borrower's contractual release rights merely to make a proposed reshape pass.

\subsection{Whole-position release and computational complexity}\label{cr:complexity}

With whole-position release, full repayment releases the entire lock; a partial repayment, if allowed, releases none. For testing the greatest cash requirement it suffices to consider loans that are either unchanged or fully repaid. A partial repayment reduces debt at an unchanged lock and cannot increase the shortfall. Thus \cref{cr:allpairs} becomes the finite family
\begin{equation}\label{cr:subsets}
\sum_{i\in J}d_i\le W+R\!\left(\sum_{i\in J}\ell_i\right)
\qquad\text{for every subset }J.
\end{equation}
For $W=0$ this reduces to the per-loan conditions $d_i\le R(\ell_i)$ by \eqref{cr:superadd}. With a shared positive reserve, testing all subsets can instead be a combinatorial problem.

\begin{samepage}
\begin{proposition}[Complexity of whole-position coverage]\label{cr:hardness}
Consider the decision problem \eqref{cr:subsets} with positive integer collateral amounts encoded in binary, positive rational debts, non-negative rational $W$, and $R(q)=q+q^2/2$ on a bounded domain containing the supply. Deciding whether the book is safe is coNP-complete. With integer locks of total $L$, a dynamic program with $O(nL)$ rational arithmetic operations decides it. Thus the hardness established here is weak; there is no claim of strong hardness or intractability for every class of curves and books.
\end{proposition}
\end{samepage}
\begin{proof}
An unsafe subset is a polynomial-size certificate: sum its integer locks and rational debts and compare them with $W+R(q)$ using exact rational arithmetic. The complement of safety is therefore in NP.

Reduce from PARTITION, whose input is a list of positive integers $w_i$ summing to $2B$, and whose question is whether a subset sums to $B$ \citep{karp1972}. The even-total restriction does not change NP-completeness, since odd-total instances are immediate no instances. Set
\[
\ell_i=w_i,\qquad d_i=(1+B)w_i,\qquad
W=\frac{B^2}{2}-\frac14,\qquad S=L=2B.
\]
All numbers have polynomial encoding length. A schedule $p(q)=1+q$ on $[0,2B+1)$ meets \cref{def:schedule}. For a retained subset of total lock $x$, the debt minus curve reserve is
\begin{equation}\label{cr:partitiondeficit}
(1+B)x-R(x)=Bx-\frac{x^2}{2}
=\frac{B^2}{2}-\frac{(x-B)^2}{2}.
\end{equation}
If $x=B$, it exceeds $W$ by $1/4$. Otherwise $x$ is an integer with $|x-B|\ge1$, so it is at most $B^2/2-1/2<W$. There is an unsafe subset exactly when the PARTITION instance is a yes instance. Hence unsafety is NP-hard and safety is coNP-hard; membership proves completeness.

For the dynamic program, let $M_j(x)$ be the largest debt of a subset of the first $j$ loans with total lock exactly $x$, or $-\infty$ if none exists. Start with $M_0(0)=0$ and all other values $-\infty$, and use
\[
M_j(x)=\max\{M_{j-1}(x),\ M_{j-1}(x-\ell_j)+d_j\},
\]
omitting the second term for $x<\ell_j$. Test $M_n(x)\le W+R(x)$ for all $0\le x\le L$. There are $O(nL)$ states, and sums and comparisons of the rational debts have polynomial bit length. This is a pseudopolynomial algorithm because the binary encoding of $L$ can be much shorter than $L$.
\end{proof}

\paragraph{Reachability of the constructed book.}
The reduction does not rely on a negative initial cash balance. It has $D=(1+B)2B=R(2B)$ and $A=W\ge0$. It can also be reached without violating instantaneous gross coverage in a model that replaces the floor loan cap by the immediate aggregate test. Buy all $2B$ tokens at curve cost, contribute $W$, and originate every lock $w_i$ with debt $w_i$. At this stage all tokens are locked and debt is $2B$. Top up each debt to $(1+B)w_i$. Total debt never exceeds $R(2B)$, cash never falls below $W$, and there are no free tokens during the top-ups. If a partition exists, fully repay the complementary loans and redeem their $B$ released tokens. The cash shortfall is $1/4$. These top-ups are deliberately outside the original floor cap; the construction is not a counterexample to \cref{thm:safety}.

The complexity comparison also has a limited meaning. Whole-position release admits fewer repayment paths than proportional release, so for the same initial book its minimum required reserve is no greater. Its subset condition can nevertheless be harder to evaluate exactly. The one-dimensional envelope solves the divisible problem because every fractional residual allocation is permitted; it is not the exact envelope for a discrete book.

\begin{proposition}[The envelope as a guard for whole-position release]\label{cr:guardbound}
For the same book and curve let $W^{\mathrm{prop}}$ be the reserve \eqref{cr:wmin} required under proportional release and $W^{\mathrm{whole}}$ the least $W$ satisfying \eqref{cr:subsets}. Then
\[
W^{\mathrm{whole}}\;\le\;W^{\mathrm{prop}}\;\le\;W^{\mathrm{whole}}+\max_i\,(d_i-\Pf\ell_i)^+ .
\]
\end{proposition}
\begin{proof}
Every whole-position configuration is a proportional one, which gives the first inequality. For the second, let $q^\ast$ maximize $H-R$ and let $j$ be the index with $Q_{j-1}\le q^\ast\le Q_j$ in the order of \cref{cr:algorithm}, so that $H(q^\ast)=B_{j-1}+\theta d_j$ and $q^\ast=Q_{j-1}+\theta\ell_j$ with $0\le\theta\le1$. The subset of the first $j-1$ loans gives $W^{\mathrm{whole}}\ge B_{j-1}-R(Q_{j-1})$. Hence
\[
W^{\mathrm{prop}}-W^{\mathrm{whole}}\le\theta d_j-\bigl[R(Q_{j-1}+\theta\ell_j)-R(Q_{j-1})\bigr]\le\theta\,(d_j-\Pf\ell_j)
\]
by \cref{lem:R}, and the right side is at most $(d_j-\Pf\ell_j)^+$.
\end{proof}

A contract with whole-position release can therefore use the sorted envelope as its guard. For a curve given by $m$ price intervals as in \cref{cr:algorithm}, the guard takes $O(n\log n+m)$ comparisons and arithmetic operations (bit complexity aside), it is safe, and it asks for more reserve than necessary by at most the largest amount by which a single loan exceeds the floor value of its lock; for an empty book that maximum is zero. The hardness result above concerns the exact requirement and is no obstacle to a design. The dynamic program can be impractical when the total collateral remains large after dividing the locks by their greatest common unit, since its table has one row per unit; the envelope guard does not depend on the magnitude of the collateral.

The constant $1$ in the bound cannot be improved in general, and it can be improved for linear prices.

\begin{corollary}[Sharpness]\label{cr:guardsharp}
(a) For every $\varepsilon\in(0,1)$ there is a schedule and a single loan with $W^{\mathrm{whole}}=0$, $(d-\Pf\ell)^+=1$ and $W^{\mathrm{prop}}=1-\varepsilon$. (b) For a linear price $p(q)=\Pf+kq$ with $k\ge0$,
\[
W^{\mathrm{prop}}-W^{\mathrm{whole}}\;\le\;\tfrac14\,\max_i\,(d_i-\Pf\ell_i)^+ ,
\]
and the factor $\tfrac14$ is attained.
\end{corollary}
\begin{proof}
(a) Let $p=1$ on $[0,1-\varepsilon)$ and $p=1+1/\varepsilon$ on $[1-\varepsilon,1)$, and take one loan with $\ell=1$, $d=2$. Then $R(1)=2=d$, so $W^{\mathrm{whole}}=0$, while $H(q)-R(q)=2q-q=q$ on $[0,1-\varepsilon]$ gives $W^{\mathrm{prop}}=1-\varepsilon$. (b) If $H-R$ is maximal at a point $Q_j$ of the partition, the gap is zero. Otherwise the maximum is at an interior point $q^\ast=Q_{j-1}+\theta\ell_j$ with $0<\theta<1$ and $k>0$, where the slope condition is $a_j=\Pf+kq^\ast$. On $[Q_{j-1},Q_j]$ the function $H-R$ is a quadratic with leading coefficient $-k/2$ and maximum at $q^\ast$, so it falls short of its maximum by $k(\theta\ell_j)^2/2$ at $Q_{j-1}$ and by $k((1-\theta)\ell_j)^2/2$ at $Q_j$. Both endpoints are whole-position configurations, so the gap is at most $\tfrac12k\ell_j^2\min\{\theta^2,(1-\theta)^2\}$. The slope condition gives $d_j-\Pf\ell_j=k\ell_jq^\ast\ge k\theta\ell_j^2$. The ratio of the two is at most $\min\{\theta^2,(1-\theta)^2\}/(2\theta)\le\tfrac14$, with equality only at $\theta=\tfrac12$ and $Q_{j-1}=0$. One loan with $\Pf=k=\ell=1$ and $d=R(1)=\tfrac32$ attains it: $W^{\mathrm{whole}}=0$, $W^{\mathrm{prop}}=\tfrac18$ and $d-\Pf\ell=\tfrac12$.
\end{proof}

Script block N15 checks the inequalities of \cref{cr:guardbound,cr:guardsharp} on 3{,}000 random books with linear prices in exact arithmetic, and evaluates the example of part~(a).

\subsection{Exact arithmetic checks}\label{cr:checks}

The standalone script \texttt{scripts/credit\_frontier.py} uses Python's \texttt{Fraction} class and has no external dependencies. Running \texttt{python3 -I scripts/credit\_frontier.py} prints the calculations and check counts; the supplied output is \texttt{outputs/credit\_frontier\_run.txt}. The tests have four distinct roles.

First, the script evaluates the numerical examples and the curve-change counterexample directly from cash flows. The irrational release in \eqref{cr:maxrelease} is bounded between two rational numbers whose reserve values bracket $3/4$; only its displayed decimal approximation uses a floating-point square root. Second, it evaluates \eqref{cr:wmin} by the sorted-envelope algorithm and independently enumerates residual collateral on a rational grid for small books. The tested curves and ratios are chosen so that every possible interval maximum lies on that grid; agreement is therefore an equality check for these fixtures, not a claim that a coarse grid generally proves safety. Third, it enumerates whole-position subsets and compares their shortfalls with both the dynamic program and independent subset-sum answers for the reduction. Finally, directed mutations replace the buffer by current aggregate coverage or an average-ratio proxy, permit its withdrawal, or retain the original reshape guards after changing the loan anchor. Each mutation admits a specified path whose cash shortfall is then computed directly.

These are deterministic finite checks of the equations and algorithms. They do not machine-check the quantified results, cover production contract arithmetic, or establish the practical desirability of the alternative contracts. Those tasks remain separate from the analytic proofs above.
